\chapter{Deux sources à la marge}\label{ch:autres}

Deux autres sources ont servi au projet Grothendieck. Elles enseignent chacune une chose
que les grands fonds n'enseignent pas.

\section{Le fonds Grothendieck--Illusie, IHÉS}

\paragraph{La source.} Un serveur de fichiers de l'IHÉS
(\url{https://repo-archives.ihes.fr/FONDS_GROTHENDIECK_ILLUSIE/}), hors du catalogue public de
l'institut, sert des notices Omeka en XML, des vues de page en JPEG d'environ
$560\times 800$ pixels et des originaux TIFF d'environ 39~Mo. Le relevé compte 198 pièces en
huit séries, environ 1\,450 pages ; une notice est illisible.

\paragraph{La leçon.} Un serveur brut n'est pas une publication. Qu'un fichier
soit accessible ne dit pas qu'il est publié, ni ce qu'on peut en faire. La correspondance
reste sous le droit de ses auteurs. Le script de catalogue attend une seconde entre deux
requêtes et ne lit pas les fichiers de reconnaissance de caractères. Le projet a d'abord
traité le fonds comme inédit, avec les transcriptions gardées en local ; il l'a ensuite
traité comme publié, pour afficher dans le volet de lecture les vues servies par l'IHÉS.
Les deux positions sont écrites dans le dépôt, et toutes deux sont à confirmer auprès des
archives de l'IHÉS. C'est exactement ce que la grille du chapitre~\ref{ch:grille} demande
de régler avant de commencer, et non pendant.

\section{Les papiers de Lasserre, BnF}

\paragraph{La source.} Les papiers trouvés chez Grothendieck à Lasserre, entrés à la BnF en
2023 par legs et don de ses enfants (NAF~29093), 43 boîtes et 22\,495 feuillets, de 1987 à
1999. Ils ne sont pas en ligne : on les consulte en salle, sur des terminaux, après une
procédure particulière auprès du département des Manuscrits.

\paragraph{La contrainte.} La méthode ne s'applique pas là où il n'y a pas d'image.
Le projet en a tiré une fiche de consultation, faite pour être imprimée et remplie au
crayon (seul instrument admis en salle), et une règle : rien de ces papiers n'est publié
sans l'accord de la BnF, qui les conserve, et des cinq héritiers, qui en détiennent les
droits d'auteur. La consultation elle-même demande une inscription au département des
Manuscrits, avec un objectif de recherche déclaré, et la signature d'une déclaration de
confidentialité ; Olivia Caramello le rappelle dans un entretien de janvier 2026 avec
Fabrice Micheau, et précise que l'Istituto Grothendieck n'a reproduit de ce fonds que des
photographies de diagrammes, avec l'autorisation des héritiers. Ce qu'on rapporte d'une
consultation se décrit dans ses mots, avec la cote et la date de la séance.

\begin{exemple}
Quand une source ne se transcrit pas, il reste d'autres façons de travailler : décrire
(la fiche de la BnF), énoncer et démontrer (les carnets de Quillen).
\end{exemple}
